\def\level{Première}  % Classe à droite
\def\course{NSI}
\def\chaptername{Représentation des entiers positifs} % Titre au centre
\def\docname{RD01}        % Identifiant à gauche

\pagestyle{cours_FP}
\makeheader{} % Affiche le bandeau

\section{Origine du binaire en informatique}

 Pour communiquer, deux êtres humains doivent parler la même langue. L'échange d'informations entre les composants d'un ordinateur repose sur le même principe.

 Dans les années 40, lorsqu'il a fallu passer de l'idée abstraite de machine de Turing (modèle mathématique) à la réalisation concrète du premier ordinateur (machine physique), la question des composants matériels à choisir s'est posée (transistor, lampe, DEL, aimant, tube, bande, \ldots{}). Pour ces composants, il est possible de réduire leur fonctionnement à 2 états exclusifs l'un de l'autre : état 0 et état 1.

 \begin{itemize}
		\item     le transistor est bloqué (état 0) ou passant (état 1),
		\item     l'interrupteur est ouvert (état 0) ou fermé (état 1),
		\item     la lampe est éteinte (état 0) ou allumée (état 1),
		\item     le ruban est trouée (état 0) ou non (état 1),	
 \end{itemize}

 	Le langage de base finalement choisie pour l'informatique, et encore utilisé aujourd'hui, est le binaire dont l'unité est le \textbf{bit} : binary digit (ou nombre binaire), pouvant prendre comme seules valeurs \textbf{0 ou 1}.


	\begin{exercice_cours}{Combien d'états différents peut-on coder avec 1 bit? 2 bits? 8 bits? $n$ bits? \dotfill}{}
	\end{exercice_cours}

	\section{Représentation binaire}

	\begin{exercice_cours}{Représentation binaire d'un entier naturel}{}

		Sur 8 bits, donner la représentation des premiers entiers naturels:

\medskip

\begin{tabular}{|c|c|c|c|c|c|c|c|c|c|c|c|c|}
	\hline
	V. décimale & V. binaire & & V. décimale & V. binaire& & V. décimale & V. binaire& & V. décimale & V. binaire\\ \hline
	0 & & & 8 & & & 16 & & & 24 &\\ \hline
	1 & & & 9 & & & 17 & & & 25 &\\ \hline 
	2 & & & 10 & & & 18 & & & 26 &\\ \hline
	3 & & & 11 & & & 19 & & & 27 &\\ \hline
	4 & & & 12 & & & 20 & & & 28 &\\ \hline
	5 & & & 13 & & & 21 & & & 29 &\\ \hline
	6 & & & 14 & & & 22 & & & 30 &\\ \hline
	7 & & & 15 & & & 23 & & & 31 &\\ \hline
	


\end{tabular}

\end{exercice_cours}

\begin{propriete}[A retenir]{}

	\begin{itemize}
		\item un entier naturel est représenté en binaire par une suite de bits (0 et 1).
		\item un entier naturel codé sur $n$ bits peut prendre $2^n$ valeurs différentes, de 0 à $2^n-1$.
		\item avec 8 bits, on peut représenter les entiers naturels de 0 à 255.
		\item le \textbf{bit de poids faible} est le bit de droite. C'est le seul multiplié par une valeur impaire. C'est donc le bit qui donnera la \textbf{parité}. Si la valeur se termine par 0, le nombre sera pair, sinon il sera impair
		\item le \textbf{bit de poids fort} est le bit de gauche.
		\item un octet est un groupe de 8 bits.
	\end{itemize}

\end{propriete}


\section{Conversion Binaire~--~Décimal: convertir $1101\,1011_2$ en décimal}

\renewcommand{\arraystretch}{1.4}
		\begin{tabular}{|p{2.6cm}|p{1.6cm}|p{1.4cm}|p{1.4cm}|p{1.4cm}|p{1.4cm}|p{1.4cm}|p{1.4cm}|p{1.4cm}|}
		\hline
		Rang du bit & 7 & 6 & 5 & 4 & 3 & 2 & 1 & 0 \\ \hline
		Puissance de 2 & $2^7=128$ & $2^6=64$ & $2^5=32$ & $2^4=16$ & $2^3=8$ & $2^2=4$ & $2^1=2$ & $2^0=1$ \\ \hline
		Bit & 1 & 1 & 0 & 1 & 1 & 0 & 1 & 1 \\ \hline
		Valeur & 128 & 64 & 0 & 16 & 8 & 0 & 2 & 1 \\ \hline
		
		\end{tabular}

\medskip

		Ainsi $1101\,1011_2 = 128+64+16+8+2+1=219_{10}$

\medskip

		\begin{exercice_cours}[Convertir les valeurs suivantes en décimal: $1001\,0100_2$, $1010_2$, $110\,1101_2$]{}
	\vspace{2cm}
	\end{exercice_cours}


	\section{Conversion Décimal~--~Binaire}

	\begin{methode}[Par la méthode des divisions successives]{}



		\medskip

		\begin{minipage}{0.74\linewidth}
					Par exemple pour convertir $44_{10}$ en binaire:

			\begin{itemize}
				\item On effectue des divisions par $2$ jusqu'à ce que le quotient soit nul.
				\item On lit ensuite les restes de la droite vers la gauche.
			
			\end{itemize}
			
			On a donc: $$44_{10}=0010 \; 1100_2$$
			
			\end{minipage}\hfill
			\begin{minipage}{0.25\linewidth}
			\begin{center}
			\includegraphics[width = 0.95\linewidth]{02.jpg}
			
			\end{center}
			
			\end{minipage}

\end{methode}

\begin{methode}[Par la méthode des soustractions successives]{}

	Lorsqu'on cherche à convertir $222_{10}$, on remarque que la plus grande puissance de 2 inférieure à 222 est $2^7=128$. Nous aurons donc besoin de 8 bits.
				
\medskip


		\begin{minipage}{0.54\linewidth}
			
			\begin{itemize}
				\item $222\geqslant 128$. On peut donc décomposer $222 = 2^7 + 94$
				\item $94 \geqslant 64$. On peut donc décomposer $94 = 2^6 + 30$
				\item $30 < 36$
				\item $30 \geqslant 16$. On peut donc décomposer $36 = 2^4 + 14$
				\item $14  \geqslant 8$. On peut donc décomposer $14 = 2^3 + 6$
				\item $6 \geqslant 4$. On peut donc décomposer $6 = 2^2 + 2$
				\item $2 \geqslant 2$. On peut donc décomposer $2 = 2^1 + 0$
				\item $0 < 1$
			\end{itemize}
			
			\end{minipage}
			\begin{minipage}{0.44\linewidth}
			\begin{center}
			\includegraphics[width = 1\linewidth]{03.jpg}
			
			\end{center}
			
			\end{minipage}
			\end{methode}

\medskip

	\begin{exercice_cours}[En utilisant les deux méthodes, convertir $234$ et $86$ sur 8 bits]{}
		\vspace{3cm}

	\end{exercice_cours}

	\pagebreak

	\section{Opérations binaires}

	\begin{propriete}[Addition en binaire]{}

		\begin{multicols}{2}

		L'addition se pratique de la même façon que dans le calcul décimal usuel, et repose sur la table d'addition binaire. 


		\columnbreak
		
		\begin{center}
		
		\begin{tabular}{|l|l|l|}
		\hline
		\rowcolor{red}+ & 0 & 1 \\ 
		\hline
		\cellcolor{red}0 & 00 & 01 \\ 
		\hline
		\cellcolor{red}1 & 01 & 10 \\ 
		\hline
		\end{tabular}
		\end{center}
		
		\end{multicols}
	\end{propriete}

	\begin{exercice_cours}{Poser et effectuer une addition en binaire: $1 \; 0101_2+1110_2$}{}
	\end{exercice_cours}

	\begin{propriete}[Multiplication en binaire]{}
		\begin{minipage}{0.8\linewidth}
			La multiplication se pratique de la même façon que dans le calcul décimal usuel, et repose sur la table de multiplication suivante :

\end{minipage}\hfill
\begin{minipage}{0.18\linewidth}
\begin{center}

\begin{tabular}{|l|l|l|}
\hline
\rowcolor{red}X & 0 & 1 \\ 
\hline
\cellcolor{red}0 & 00 & 00 \\ 
\hline
\cellcolor{red}1 & 00 & 01 \\ 
\hline
\end{tabular}
\end{center}
\end{minipage}
	\end{propriete}
	
	\begin{exercice_cours}[Poser et effectuer une multiplication en binaire: $1 \; 0101_2 \times 100_2$]{}
	\vspace{2cm}
	\end{exercice_cours}

	\begin{propriete}[A retenir]{}
		\begin{itemize}
			\item \textbf{Multiplier} une valeur décimale par 2 revient à décaler tous les bits d'un cran vers la gauche et à ajouter un 0 sur le bit de poids faible.
			\item \textbf{Diviser} une valeur décimale par 2 revient à décaler tous les bits d'un cran vers la droite. Le bit de poids faible est perdu. Cela revient à récupérer le quotient de la division par 2.
			\item \textbf{Multiplier} une valeur décimale par $2^n$ revient à décaler tous les bits de $n$ crans vers la gauche et à ajouter des 0 sur les bits de poids faible.
			\item \textbf{Diviser} une valeur décimale par $2^n$ revient à décaler tous les bits de $n$ crans vers la droite. Les bit de poids faible sont perdus. Cela revient à récupérer le quotient de la division par $2^n$.
		\end{itemize}
	\end{propriete}


	
		\section{L'Hexadécimal}


	\begin{propriete}[Représentation d'un entier en hexadécimal]{}
		Si le binaire est une base comprenant 2 éléments, l'hexadécimal est une base comprenant 16 éléments. Il est donc également possible de coder les valeurs décimales suivantes en hexadécimal.

Puisqu'il nous faut 16 symboles. On utilise les 10 premiers chiffres et on rajouter les symboles $\{A, B, C, D, E, F\}$

\begin{center}
	\includegraphics[scale=0.8]{04.jpg}
	\end{center}
	\end{propriete}

	
	\begin{propriete}[Conversion hexadécimale]{}
Les méthodes utilisées sont les memes que précédemment en remplaçant la valeur 2 par 16
	\end{propriete}

	\begin{exercice_cours}[Convertir en décimal: $A3_{16}$, $BCF_{16}$, $42_{16}$ ]{}
	\vspace{2cm}
	\end{exercice_cours}

	\begin{exercice_cours}[Convertir en hexadécimal: $364_{10}$, $23_{10}$, $44_{10}$ ]{}
	\vspace{2cm}
	\end{exercice_cours}

	\begin{propriete}[Conversion binaire -- hexadécimal]{}
		Dans le cas de conversion binaire - hexadécimal (et inversement) il n'est pas utile de repasser par le décimal. En effet un symbole hexadécimal peut directement etre codé par 4 bits.

		Pour convertir des entiers naturels de la représentation binaire à la représentation hexadécimale: Il suffit de regrouper les bits par quatre et de leur associer leur chiffre hexadécimal.
		
		$0110\,1000_2 = 68_{16}$
	\end{propriete}


	\begin{minipage}{0.49\linewidth}
		
	\begin{exercice_cours}[Convertir en hexadécimal:]{}
		\begin{itemize}
			\item $1101\,1010\,0011_2$
			\item $0011\,1111\,1001_2$
		\end{itemize}
	\end{exercice_cours}	\end{minipage}
\hfill
	\begin{minipage}{0.49\linewidth}
		\begin{exercice_cours}[Convertir en binaire:]{}
		\begin{itemize}
			\item $A3E_{16}$
			\item $8DC_{16}$
		\end{itemize}
	\end{exercice_cours}
	\end{minipage}
	
\begin{propriete}[Conversion en base $b$]{}
	Dans le même ordre d'idée il est donc possible de faire des conversion depuis n'importe quelle base $b$.

	\bigskip
	
	Pour tout entier naturel $b$ non nul appelé base, tout nombre entier naturel $n$ peut se décomposer sous la forme:
	
	$$n=\sum_{k = 0}^{N} a_kb^k=a_Nb^N+ \ldots + a_1b^1+a_0b^0$$
	
	L'écriture en base $b$ de $n$ sera alors $\overline{a_Na_{N-1} \ldots a_1a_0}$
	
\end{propriete}

\begin{methode}[Conversion en python]


\begin{CodePythontexAlt}[Largeur=1\linewidth,Centre,PremLigne=1, lineNumbers=false]{}
a_bin = 0b1101	# le 0b veut dire que la valeur est écrite en binaire
a_dec = int(a_bin)	# conversion en décimal - valeur entière égale à 13
a_hex = hex(a_dec)	# conversion en hexadécimal - valeur égale à 0xd

b_dec = 44
b_bin = bin(b_dec)	# conversion en binaire - valeur égale à 0b101100
b_hex = hex(b_dec)	# conversion en hexadécimal - valeur égale à 0x2c

# de manière générale, pour convertir un entier n en base b, on peut utiliser:
a_dec = int('1101', 2)	 # conversion binaire en décimal égale à 13
a_hex = hex(int('1101', 2))	# conversion binaire en hexadécimal égale à 0xd
a_bin = bin(int('A3', 16))	# conversion hexadécimal en binaire égale à 0b10100011
\end{CodePythontexAlt}
\end{methode}
	\pagebreak
